\scanpage{32}
% L1–5
\deleted{\unsure{Réciproquement|Prouvons|Démonstration} :}

$a\Rightarrow b\Rightarrow c$ trivial,

$c\Rightarrow a$ car $\unsure{N\mapsto|\mathstrut N\simeq}\operatorname{Ext}^{n}(N,M)$ sera exact, donc $N\longrightarrow\operatorname{Ext}^{n+1}(N,M)$ sera exact à gauche, donc $\operatorname{Ext}^{n+1}(k,M)\to H^{n+1}(M)$ sera injectif, or $H^{n+1}(M)=0$.

% L6–10
\uline{Corollaire 1} Sous ces conditions, \inserted{on a $H^n(M)\simeq D(A)^m$ pour certain $m$,} \deleted{les} idéaux premiers \deleted{et \unsure{ce} sont les idéaux pr. associés \unsure{$M$},} \deleted{de} \inserted{à $0$ dans} $A$ sont tous de dim $n$, et on a pour tout \deleted{\unsure{idéal}} tel $\mathfrak p$
\[
 \lg_{\mathfrak p}(M)=m\lg_{\mathfrak p}(A)
 \qquad \begin{aligned}\operatorname{Ext}^{n}(k,M)&\simeq k^{m}\\
 \operatorname{Ext}^{i}(k,M)&=0\ \text{si } i>n\end{aligned}
\]

% L11–12
\deleted{\unsure{Corollaire|Cor.} 2 \unsure{Si} $0<\lg_{\mathfrak p}M\leq\lg_{\mathfrak p}A$, les conditions a) b) c) \unsure{qui précèdent} équivalent à b') $\unsure{\operatorname{Ext}^{n}(N,M)|\operatorname{Ext}^{n}(M,M)}$ \unsure{est} c') $H^{n}(M)$ est dualisant.}

% margin note written sideways
\margin{\uline{Cor 2} Les conditions précédentes jointes à $0<\lg_{\mathfrak p}(M)\leq\lg_{\mathfrak p}(A)$ équivalent aux conditions \deleted{\unsure{équiv.|équivalentes}} suivantes compte tenu de $\operatorname{Ext}^{n-1}(k,M)=0$.\newline a) $\operatorname{Ext}^{n}(k,M)=k$, $\operatorname{Ext}^{n+1}(k,M)=0$\newline b) $N\longrightarrow\operatorname{Ext}^{n}(N,M)$ est dualisant\newline c) $H^{n}(M)$ est dualisant}

% L13–19
\uline{Remarque} \deleted{On ignore} \inserted{[\unsure{à reporter|reporter|injectif} à la fin du N° \unsure{+}]} Il semble que ces conditions impliquent que \deleted{\unsure{$B$}} $A$ est Coh.-Mac, et \deleted{\unsure{$B=\operatorname{Hom}$}} $M\simeq\Omega^{m}$, où $\Omega$ est un module fondamental pour $A$, \inserted{[et que l'hyp. $\operatorname{Ext}^{n-1}(k,M)=0$ soit $=$ redondante dans b')]} Pour le prouver, procédons par réc. sur $n$, il suffirait de prouver que pour $n>1$, l'annulateur de $\mathfrak m$ dans $M$ est nul. Voici un résultat partiel dans ce sens

% margin note written sideways
\margin{\uline{Cor 3} Les conditions du Cor 2 sont entraînées par l'une quelcon. des cond. suivantes\newline d) $\operatorname{Ext}^{n}(k,M)\simeq k$, $M$ fidèle\newline \unsure{d')bis)|d bis)} $\operatorname{Ext}^{\unsure{r|n}}(\unsure{A/\mathfrak n|A/\mathfrak m|A/\mathfrak p|A/\mathfrak N},M)$ est \deleted{\unsure{de|co}} \unsure{comonogène|monogène}, $M$ fidèle\newline \uline{Prouver que} \deleted{\unsure{$M$ soit}} ces idéaux pr. associés à $M$ sont tous de dim $n$. \uline{[\unsure{question}]}}

% L20–24
\uline{Proposition 14} \deleted{Sous les conditions précédentes,} \inserted{Supposons $m=1$, \unsure{ou} $A$ un quotient de régulier.} supposons de plus $M$. Coh. Mac. Alors les conditions a) b) c) \deleted{\unsure{sont encore équivalentes|sont chacune équivalente|sont devenues équivalentes}} à :

$A$ est Coh. Mac. et $M\simeq\Omega_A^{m}$, où $\Omega_A$ est un module fondamental pour $A$

% L25–26
C'est trivial si $n=0$, procédons par réc. sur $n$. Voir les notes jaunes - - - - -

% archive sheet number: 31
